\scanpage{29}
% Pencil archive foliation: 28.
$n$ assez grand, sous la forme
$M_n(n)$, avec $M_n\in\operatorname{Ob}\mathcal{M}^+(X)$, [pour $n$ fixé,
$M_n$ est bien déterminé par $M$ à isom unique
près, et
\[
(M_n(n)\simeq M_{n+1}(n+1))\Longleftrightarrow(M_{n+1}\simeq M_n(-1)),
\]
on retrouve la description de $\mathcal{M}(X)$ en
termes de \unsure{pseudo-}$\underset{\longrightarrow}{\lim}$.]
Les foncteurs
$T_\ell$ s'étendent à $\mathcal{M}(X)$, de façon unique, de façon
à rester compatible avec $\otimes$.

7) Dans $\mathcal{M}(X)$, on a aussi des foncteurs
$\underline{\operatorname{Hom}}$, liés à $\otimes$ par la formule
habituelle
\[
\operatorname{Hom}(P\otimes Q,R)\simeq
\operatorname{Hom}(P,\underline{\operatorname{Hom}}(Q,R))
\simeq\operatorname{Hom}(Q,\underline{\operatorname{Hom}}(P,R)),
\]
et \deleted{satisfaisant}\inserted{qui s'étendent en}
\[
R\operatorname{Hom}(\unsure{P|R},Q)
\qquad P,Q\in\operatorname{Ob}D^bM(X)
\]
satisfaisant à la relation analogue
relativement à $\underline{\otimes}$.
\uuline{La formation des $\underline{\operatorname{Hom}}$
\inserted{et $R\underline{\operatorname{Hom}}$}
est compatible avec les} $T_\ell$.

\deleted{NB} [NB. On retrouve la formation
des $f^!$ \ldots]
